\documentclass[UTF8,a4paper,12pt]{ctexart}

\usepackage{geometry}
\usepackage{hyperref}
\usepackage{xcolor}
\usepackage{listings}
\usepackage{enumitem}
\usepackage{booktabs}
\usepackage{amsmath}
\usepackage{amssymb}

\geometry{left=2.5cm,right=2.5cm,top=2.5cm,bottom=2.5cm}
\hypersetup{
  colorlinks=true,
  linkcolor=blue,
  urlcolor=blue
}

\lstdefinestyle{code}{
  basicstyle=\ttfamily\small,
  breaklines=true,
  frame=single,
  columns=fullflexible,
  backgroundcolor=\color{gray!6},
  rulecolor=\color{gray!40}
}
\lstset{style=code}

\title{脑电运动解码入门环境配置与使用说明}
\author{}
\date{\today}

\begin{document}

\maketitle

\section{环境用途}

本环境用于入门级脑电信息解码实验，尤其适合先做运动想象（Motor Imagery, MI）任务，例如左手与右手运动想象二分类。环境已将主要软件包、数据集目录、缓存目录配置到 F 盘，避免公开数据集和 Python 包大量占用 C 盘。

推荐的入门实验流程如下：

\begin{enumerate}[label=\arabic*.]
  \item 使用公开运动想象 EEG 数据集。
  \item 用 MNE 做读取、滤波、分段和可视化。
  \item 用 CSP、LDA、SVM 或黎曼几何方法做分类。
  \item 用交叉验证评估准确率。
  \item 后续再升级到深度学习方法，例如 EEGNet、Braindecode。
\end{enumerate}

\section{目录结构}

当前配置的主要目录如下：

\begin{table}[h]
\centering
\begin{tabular}{ll}
\toprule
路径 & 作用 \\
\midrule
\texttt{F:\textbackslash Miniforge3} & Conda/Miniforge 管理器安装位置 \\
\texttt{F:\textbackslash eeg\_bci\textbackslash envs\textbackslash eeg-mi} & 脑电实验 Python 环境 \\
\texttt{F:\textbackslash eeg\_bci\textbackslash data} & EEG 数据集根目录 \\
\texttt{F:\textbackslash eeg\_bci\textbackslash data\textbackslash mne} & MNE 下载数据目录 \\
\texttt{F:\textbackslash eeg\_bci\textbackslash data\textbackslash moabb} & MOABB 下载数据目录 \\
\texttt{F:\textbackslash eeg\_bci\textbackslash projects} & 建议存放代码、Notebook 和实验结果 \\
\texttt{F:\textbackslash eeg\_bci\textbackslash pkgs} & Conda/Pip 包缓存目录 \\
\texttt{F:\textbackslash eeg\_bci\textbackslash tmp} & 临时文件目录 \\
\bottomrule
\end{tabular}
\end{table}

\section{已安装的软件包}

环境名称为 \texttt{eeg-mi}，Python 版本为 \texttt{3.11.15}。已验证的软件包包括：

\begin{table}[h]
\centering
\begin{tabular}{lll}
\toprule
软件包 & 版本 & 用途 \\
\midrule
\texttt{mne} & \texttt{1.12.1} & EEG/MEG 数据读取、预处理、滤波、分段、可视化 \\
\texttt{moabb} & \texttt{1.5.0} & 公开 BCI 数据集下载与统一评测 \\
\texttt{scikit-learn} & \texttt{1.9.0} & 机器学习分类器、交叉验证、指标评估 \\
\texttt{pyriemann} & \texttt{0.12} & 基于黎曼几何的 EEG 分类方法 \\
\texttt{jupyterlab} & 已安装 & 交互式编程和实验记录 \\
\bottomrule
\end{tabular}
\end{table}

\section{环境变量}

以下路径已经写入系统环境变量和 Conda 环境变量。重新打开终端后生效。

\begin{lstlisting}[language={},caption={已配置的环境变量}]
MNE_DATA=F:\eeg_bci\data\mne
MOABB_DATA=F:\eeg_bci\data\moabb
JOBLIB_TEMP_FOLDER=F:\eeg_bci\tmp
MPLCONFIGDIR=F:\eeg_bci\tmp\matplotlib
PIP_CACHE_DIR=F:\eeg_bci\pkgs\pip-cache
\end{lstlisting}

这些变量的作用是让 MNE、MOABB、Matplotlib、Joblib 和 Pip 优先使用 F 盘目录。

\section{启动 JupyterLab}

平时建议使用 JupyterLab 做脑电实验。打开 PowerShell，运行：

\begin{lstlisting}[language={},caption={启动 JupyterLab}]
F:\Miniforge3\Scripts\conda.exe run -p F:\eeg_bci\envs\eeg-mi jupyter lab --notebook-dir F:\eeg_bci\projects
\end{lstlisting}

运行后终端会显示一个类似下面的地址：

\begin{lstlisting}[language={}]
http://localhost:8888/lab?token=...
\end{lstlisting}

在浏览器中打开该地址即可使用。Notebook 和实验代码建议保存在：

\begin{lstlisting}[language={}]
F:\eeg_bci\projects
\end{lstlisting}

\section{运行 Python 脚本}

如果不使用 JupyterLab，也可以直接运行某个 Python 文件：

\begin{lstlisting}[language={},caption={运行 Python 文件}]
F:\Miniforge3\Scripts\conda.exe run -p F:\eeg_bci\envs\eeg-mi python F:\eeg_bci\projects\test.py
\end{lstlisting}

\section{环境测试代码}

在 JupyterLab 中新建 Notebook，运行下面代码。如果能够输出版本号，说明环境可用。

\begin{lstlisting}[language=Python,caption={验证核心软件包}]
import os
import mne
import moabb
import sklearn
import pyriemann

print("mne:", mne.__version__)
print("moabb:", moabb.__version__)
print("sklearn:", sklearn.__version__)
print("pyriemann:", pyriemann.__version__)
print("MNE_DATA:", os.environ.get("MNE_DATA"))
print("MOABB_DATA:", os.environ.get("MOABB_DATA"))
\end{lstlisting}

\section{第一个运动想象数据集测试}

可以先测试 MOABB 是否能识别公开数据集。例如使用经典的 \texttt{BNCI2014\_001} 数据集：

\begin{lstlisting}[language=Python,caption={查看公开运动想象数据集信息}]
from moabb.datasets import BNCI2014_001

dataset = BNCI2014_001()

print("Subjects:", dataset.subject_list[:5])
print("Events:", dataset.event_id)
\end{lstlisting}

该数据集包含左手、右手、脚、舌头等运动想象任务。入门时建议先只做左手与右手二分类。

\section{推荐入门实验路线}

建议先按下面顺序推进：

\begin{enumerate}[label=\arabic*.]
  \item 跑通 \texttt{BNCI2014\_001} 数据集读取。
  \item 只选择左手和右手两类。
  \item 对 EEG 进行 \texttt{8--30 Hz} 滤波。
  \item 取提示后约 \texttt{0.5--4 s} 的时间窗。
  \item 使用 CSP 提取空间特征。
  \item 使用 LDA 或 SVM 分类。
  \item 使用交叉验证报告准确率和平衡准确率。
\end{enumerate}

不要一开始就直接做深度学习。先用 CSP + LDA 建立一个稳定基线，后续再比较 EEGNet、ShallowConvNet 等方法。

\section{后续升级}

如果以后要做深度学习，可以在当前环境中追加安装 PyTorch 和 Braindecode：

\begin{lstlisting}[language={},caption={升级到深度学习环境}]
F:\Miniforge3\Scripts\conda.exe run -p F:\eeg_bci\envs\eeg-mi pip install -i https://pypi.tuna.tsinghua.edu.cn/simple torch braindecode
\end{lstlisting}

如果后续需要 GPU 版 PyTorch，应根据电脑显卡和 CUDA 版本选择对应安装命令。没有 GPU 时也可以先使用 CPU 版，入门实验和传统机器学习方法不依赖 GPU。

\section{内存与磁盘建议}

当前入门环境大约占用如下空间：

\begin{table}[h]
\centering
\begin{tabular}{ll}
\toprule
路径 & 大约占用 \\
\midrule
\texttt{F:\textbackslash Miniforge3} & \texttt{0.62 GB} \\
\texttt{F:\textbackslash eeg\_bci\textbackslash envs\textbackslash eeg-mi} & \texttt{1.97 GB} \\
\texttt{F:\textbackslash eeg\_bci\textbackslash pkgs} & \texttt{2.33 GB} \\
\texttt{F:\textbackslash eeg\_bci\textbackslash data} & 当前约 \texttt{0 GB} \\
\bottomrule
\end{tabular}
\end{table}

建议配置：

\begin{itemize}
  \item 最低内存：\texttt{8 GB}。
  \item 推荐内存：\texttt{16 GB} 或以上。
  \item 初期硬盘空间：至少预留 \texttt{10--20 GB}。
  \item 多数据集实验：建议预留 \texttt{50 GB} 或以上。
\end{itemize}

降低内存占用的做法：

\begin{lstlisting}[language=Python,caption={降低 EEG 处理内存占用的常用做法}]
raw = mne.io.read_raw_fif(path, preload=False)
raw.pick_types(eeg=True)
raw.filter(8, 30)
raw.resample(128)
\end{lstlisting}

做 MOABB 实验时，也建议先单个被试、单个数据集跑通，不要一开始批量跑所有被试和所有数据集。

\section{常见问题}

\subsection{是否容易升级？}

容易。当前环境已经包含 MNE、MOABB、Scikit-learn 和 Pyriemann，适合传统运动想象解码。后续只需要在同一环境中追加 PyTorch、Braindecode 或其他库即可。

\subsection{会不会占 C 盘？}

主要环境、包缓存和数据集目录都已经放在 F 盘。C 盘可能仍会保存少量 Windows 用户配置文件，但不会是主要数据和包存放位置。

\subsection{为什么不用系统自带 Python？}

系统 Python 当前是 \texttt{3.12.3}。脑电分析工具链里有些包对版本和依赖比较敏感，使用独立 Conda 环境可以减少冲突，也方便以后升级、删除或复现实验。

\subsection{以后怎么删除这个环境？}

如果以后确认不需要该环境，可以执行：

\begin{lstlisting}[language={},caption={删除 eeg-mi 环境}]
F:\Miniforge3\Scripts\conda.exe env remove -p F:\eeg_bci\envs\eeg-mi
\end{lstlisting}

删除前应先确认自己的代码和数据是否已经备份。一般情况下不建议删除下面目录中的项目文件：

\begin{lstlisting}[language={}]
F:\eeg_bci\projects
\end{lstlisting}

\section{入门实验操作步骤与原理}

本节给出第一个完整实验：使用公开数据集 \texttt{BNCI2014\_001}，完成左手与右手运动想象二分类。目标是先跑出一个可靠的传统机器学习基线：

\begin{lstlisting}[language={}]
EEG数据 -> 8-30 Hz滤波 -> 0.5-4.0 s分段 -> CSP特征 -> LDA分类 -> 交叉验证准确率
\end{lstlisting}

\subsection{启动 JupyterLab}

打开 PowerShell，运行：

\begin{lstlisting}[language={},caption={启动入门实验 Notebook}]
F:\Miniforge3\Scripts\conda.exe run -p F:\eeg_bci\envs\eeg-mi jupyter lab --notebook-dir F:\eeg_bci\projects
\end{lstlisting}

在浏览器中打开终端给出的 JupyterLab 地址，新建 Notebook，建议命名为：

\begin{lstlisting}[language={}]
motor_imagery_baseline.ipynb
\end{lstlisting}

\subsection{第一步：测试环境}

在第一个 Notebook 单元中运行：

\begin{lstlisting}[language=Python,caption={测试脑电实验环境}]
import os
import mne
import moabb
import sklearn
import pyriemann

print("mne:", mne.__version__)
print("moabb:", moabb.__version__)
print("sklearn:", sklearn.__version__)
print("pyriemann:", pyriemann.__version__)
print("MNE_DATA:", os.environ.get("MNE_DATA"))
print("MOABB_DATA:", os.environ.get("MOABB_DATA"))
\end{lstlisting}

如果能够输出版本号，并且 \texttt{MNE\_DATA}、\texttt{MOABB\_DATA} 指向 F 盘，说明环境正常。

\subsection{第二步：查看公开数据集}

\begin{lstlisting}[language=Python,caption={查看 \texttt{BNCI2014\_001} 数据集信息}]
from moabb.datasets import BNCI2014_001

dataset = BNCI2014_001()

print("被试列表:", dataset.subject_list)
print("任务标签:", dataset.event_id)
\end{lstlisting}

应能看到类似结果：

\begin{lstlisting}[language={}]
{'left_hand': 1, 'right_hand': 2, 'feet': 3, 'tongue': 4}
\end{lstlisting}

入门实验先只使用 \texttt{left\_hand} 和 \texttt{right\_hand} 两类。

\subsection{第三步：读取一个被试的数据}

第一次运行下面代码时会下载数据。由于前面已经设置了环境变量，数据会保存到 F 盘的数据目录中。

\begin{lstlisting}[language=Python,caption={读取一个被试的左右手运动想象数据}]
from moabb.paradigms import MotorImagery

paradigm = MotorImagery(
    events=["left_hand", "right_hand"],
    fmin=8,
    fmax=30,
    tmin=0.5,
    tmax=4.0,
    resample=128
)

X, y, metadata = paradigm.get_data(
    dataset=dataset,
    subjects=[1]
)

print("X shape:", X.shape)
print("y shape:", y.shape)
print("标签:", set(y))
print(metadata.head())
\end{lstlisting}

\texttt{X} 的形状通常是：

\begin{lstlisting}[language={}]
试次数 x 通道数 x 时间点数
\end{lstlisting}

例如 \texttt{(144, 22, 448)} 表示：144 个运动想象片段、22 个 EEG 通道、每段 448 个采样点。

\subsection{第四步：CSP + LDA 分类}

\begin{lstlisting}[language=Python,caption={CSP + LDA 交叉验证分类}]
from sklearn.pipeline import make_pipeline
from sklearn.model_selection import StratifiedKFold, cross_val_score
from sklearn.discriminant_analysis import LinearDiscriminantAnalysis
from mne.decoding import CSP

clf = make_pipeline(
    CSP(n_components=6, reg=None, log=True, norm_trace=False),
    LinearDiscriminantAnalysis()
)

cv = StratifiedKFold(
    n_splits=5,
    shuffle=True,
    random_state=42
)

scores = cross_val_score(
    clf,
    X,
    y,
    cv=cv,
    scoring="balanced_accuracy"
)

print("每折准确率:", scores)
print("平均准确率:", scores.mean())
\end{lstlisting}

如果平均准确率明显高于 \texttt{0.5}，说明模型已经学到了一定的左右手运动想象差异。

\subsection{第五步：绘制混淆矩阵}

\begin{lstlisting}[language=Python,caption={绘制混淆矩阵}]
from sklearn.model_selection import cross_val_predict
from sklearn.metrics import ConfusionMatrixDisplay
import matplotlib.pyplot as plt

y_pred = cross_val_predict(clf, X, y, cv=cv)

ConfusionMatrixDisplay.from_predictions(y, y_pred)
plt.title("Left hand vs Right hand")
plt.show()
\end{lstlisting}

混淆矩阵可以显示左手和右手分别被正确分类、错误分类的次数。它比单独看一个平均准确率更直观。

\subsection{实验原理}

运动想象解码不是直接判断原始 EEG 波形是否相似，而是利用运动想象在特定脑区和特定频段引起的节律能量变化。

左右手运动想象主要涉及感觉运动皮层：

\begin{itemize}
  \item 左手运动想象通常更明显影响右侧运动皮层。
  \item 右手运动想象通常更明显影响左侧运动皮层。
  \item 常关注的节律包括 \texttt{8--13 Hz} 的 Mu 节律和 \texttt{13--30 Hz} 的 Beta 节律。
\end{itemize}

因此入门实验中使用 \texttt{8--30 Hz} 滤波，目的是保留和运动想象最相关的频段，同时减少低频漂移和高频噪声。

时间窗设置为 \texttt{0.5--4.0 s}，表示从提示出现后 0.5 秒开始截取到 4 秒。这样可以避开刚看到提示时的瞬时视觉反应和准备阶段，使分类特征更集中在运动想象过程。

\subsection{CSP 的含义}

CSP 是 Common Spatial Pattern，中文可理解为公共空间模式。它的目标是学习一组空间滤波器，使两类任务的方差差异尽可能大。

直观理解：

\begin{lstlisting}[language={}]
找到一些 EEG 通道的加权组合，
让左手想象时能量较强、右手想象时能量较弱，
或者让右手想象时能量较强、左手想象时能量较弱。
\end{lstlisting}

EEG 单个通道信噪比较低，所以 CSP 不只依赖某一个电极，而是自动学习多个电极的组合。CSP 输出的是每个 trial 的空间模式能量特征，而不是原始波形。

\subsection{LDA 的含义}

LDA 是 Linear Discriminant Analysis，即线性判别分析。它会在 CSP 提取出的特征空间中寻找一条线性分类边界，把左手和右手尽量分开。

LDA 适合作为入门基线，原因是：

\begin{itemize}
  \item 对小样本 EEG 数据比较友好。
  \item 计算量小，结果稳定。
  \item 与 CSP 搭配是运动想象 BCI 中非常经典的传统方法。
\end{itemize}

\subsection{为什么要用交叉验证}

不能只在训练数据上报告准确率，因为这样会高估模型性能。交叉验证会把数据分成若干份，例如 5 份：

\begin{lstlisting}[language={}]
每次用 4 份训练，1 份测试；
重复 5 次；
最后报告平均测试准确率。
\end{lstlisting}

代码中使用 \texttt{make\_pipeline(CSP, LDA)} 的原因是：每一折交叉验证中，CSP 只会在训练集上学习空间滤波器，然后再应用到测试集。这样可以避免数据泄漏。

\subsection{完成本实验后的扩展方向}

跑通第一个 baseline 后，可以按下面顺序扩展：

\begin{enumerate}[label=\arabic*.]
  \item 更换不同被试，观察每个被试的准确率差异。
  \item 调整频段，例如 \texttt{8--13 Hz}、\texttt{13--30 Hz}、\texttt{8--35 Hz}。
  \item 调整时间窗，例如 \texttt{1.0--4.0 s} 或 \texttt{0.5--3.0 s}。
  \item 尝试 FBCSP，即多个频段分别做 CSP 后再分类。
  \item 使用 PyRiemann 方法，与 CSP + LDA 对比。
  \item 最后再升级到 Braindecode、EEGNet 等深度学习方法。
\end{enumerate}

\section{遇到问题和解决方法}

本节记录配置和使用过程中已经遇到的问题。遇到类似现象时，优先按本节排查。

\subsection{\texttt{Select Kernel} 中没有 \texttt{eeg-mi}}

如果 JupyterLab 的 \texttt{Select Kernel} 列表中没有看到 \texttt{eeg-mi} 或 \texttt{Python (eeg-mi)}，不一定代表环境不能用。先在 Notebook 中运行：

\begin{lstlisting}[language=Python,caption={确认当前 Notebook 使用的 Python 解释器}]
import sys
print(sys.executable)
\end{lstlisting}

如果输出类似：

\begin{lstlisting}[language={}]
F:\eeg_bci\envs\eeg-mi\python.exe
\end{lstlisting}

说明当前 Notebook 实际上已经在使用正确的脑电环境。即使 Kernel 名字显示为 \texttt{Python 3}，也可以继续实验。

进一步验证核心包：

\begin{lstlisting}[language=Python,caption={确认脑电包可用}]
import mne
import moabb

print(mne.__version__)
print(moabb.__version__)
\end{lstlisting}

如果可以正常输出版本号，说明环境可用。

\subsection{手动注册 \texttt{Python (eeg-mi)} Kernel}

如果希望 \texttt{Select Kernel} 中明确显示 \texttt{Python (eeg-mi)}，可以关闭当前 JupyterLab 后，在 PowerShell 中运行下面命令。

先创建 Jupyter 相关目录：

\begin{lstlisting}[language={},caption={创建 Jupyter 配置目录}]
New-Item -ItemType Directory -Force -Path "F:\eeg_bci\jupyter\config"
New-Item -ItemType Directory -Force -Path "F:\eeg_bci\jupyter\data"
New-Item -ItemType Directory -Force -Path "F:\eeg_bci\jupyter\runtime"
New-Item -ItemType Directory -Force -Path "F:\eeg_bci\tmp"
\end{lstlisting}

然后设置当前 PowerShell 会话的 Jupyter 路径：

\begin{lstlisting}[language={},caption={设置 Jupyter 和临时目录到 F 盘}]
$env:JUPYTER_CONFIG_DIR="F:\eeg_bci\jupyter\config"
$env:JUPYTER_DATA_DIR="F:\eeg_bci\jupyter\data"
$env:JUPYTER_RUNTIME_DIR="F:\eeg_bci\jupyter\runtime"
$env:TMP="F:\eeg_bci\tmp"
$env:TEMP="F:\eeg_bci\tmp"
\end{lstlisting}

注册内核：

\begin{lstlisting}[language={},caption={注册 eeg-mi 内核}]
F:\eeg_bci\envs\eeg-mi\python.exe -m ipykernel install --sys-prefix --name eeg-mi --display-name "Python (eeg-mi)"
\end{lstlisting}

成功后通常会出现类似输出：

\begin{lstlisting}[language={}]
Installed kernelspec eeg-mi in ...
\end{lstlisting}

之后重新启动 JupyterLab，再选择 \texttt{Python (eeg-mi)}。

\subsection{\texttt{New-Item}: 路径中具有非法字符}

如果 PowerShell 报：

\begin{lstlisting}[language={}]
New-Item : 路径中具有非法字符。
\end{lstlisting}

通常是复制命令时路径中混入了换行或隐藏字符。例如路径被复制成了类似：

\begin{lstlisting}[language={}]
F:\
  \eeg_bci\jupyter
\end{lstlisting}

这种路径是非法的。解决方法是：不要复制很长的一整行命令，改成一行一个目录，手动确认路径中间没有换行：

\begin{lstlisting}[language={},caption={逐行创建目录}]
New-Item -ItemType Directory -Force -Path "F:\eeg_bci\jupyter\config"
New-Item -ItemType Directory -Force -Path "F:\eeg_bci\jupyter\data"
New-Item -ItemType Directory -Force -Path "F:\eeg_bci\jupyter\runtime"
New-Item -ItemType Directory -Force -Path "F:\eeg_bci\tmp"
\end{lstlisting}

也可以使用更短的写法：

\begin{lstlisting}[language={},caption={使用 mkdir 创建目录}]
cd F:\
mkdir eeg_bci
mkdir eeg_bci\jupyter
mkdir eeg_bci\jupyter\config
mkdir eeg_bci\jupyter\data
mkdir eeg_bci\jupyter\runtime
mkdir eeg_bci\tmp
\end{lstlisting}

如果提示目录已经存在，可以忽略。

\subsection{\texttt{New-Item}: 找不到 \texttt{LiteralPath} 参数}

有些 PowerShell 环境中 \texttt{New-Item} 不支持 \texttt{-LiteralPath} 参数。如果出现：

\begin{lstlisting}[language={}]
New-Item : 找不到与参数名称“LiteralPath”匹配的参数。
\end{lstlisting}

把 \texttt{-LiteralPath} 改成 \texttt{-Path} 即可：

\begin{lstlisting}[language={},caption={使用 Path 参数}]
New-Item -ItemType Directory -Force -Path "F:\eeg_bci\jupyter\config"
\end{lstlisting}

\subsection{\texttt{IProgress not found} 警告}

在 Notebook 中导入 MOABB 或运行数据下载时，可能出现：

\begin{lstlisting}[language={}]
TqdmWarning: IProgress not found. Please update jupyter and ipywidgets.
\end{lstlisting}

这不是错误。它表示 Jupyter 中的进度条组件 \texttt{ipywidgets} 没有安装或没有启用，影响的是下载或运行时的进度条显示，不影响 MNE、MOABB、CSP、LDA 等核心实验代码运行。

如果环境测试输出如下内容，说明环境本身正常：

\begin{lstlisting}[language={}]
mne: 1.12.1
moabb: 1.5.0
sklearn: 1.9.0
pyriemann: 0.12
MNE_DATA: F:\eeg_bci\data\mne
MOABB_DATA: F:\eeg_bci\data\moabb
\end{lstlisting}

若希望消除该警告，可以关闭 JupyterLab 后，在 PowerShell 中运行：

\begin{lstlisting}[language={},caption={安装 ipywidgets}]
F:\Miniforge3\Scripts\conda.exe install -p F:\eeg_bci\envs\eeg-mi -c conda-forge ipywidgets -y
\end{lstlisting}

安装完成后重新启动 JupyterLab。

\subsection{\texttt{paradigm.get\_data()} 报错的修复方案}

如果运行下面代码时在 \texttt{paradigm.get\_data()} 附近报错：

\begin{lstlisting}[language=Python]
X, y, metadata = paradigm.get_data(
    dataset=dataset,
    subjects=[1]
)
\end{lstlisting}

常见原因不是分类算法错误，而是第一次下载或读取 BNCI 数据集时，MNE/MOABB 尝试写入 C 盘用户配置目录，例如：

\begin{lstlisting}[language={}]
C:\Users\LENOVO\.mne
\end{lstlisting}

如果该目录权限异常，就会出现 \texttt{PermissionError}。此外，MOABB 1.5 在 Windows 下还可能把盘符 \texttt{F:} 中的冒号错误清洗成 \texttt{F-}，导致数据路径跑到当前工作目录下面。

推荐在 Notebook 最前面使用下面的完整修复单元。运行后，建议重启 Kernel，并从第一个单元开始依次运行。

\begin{lstlisting}[language=Python,caption={修复 MNE/MOABB 数据路径与 Windows 盘符问题}]
import os
from pathlib import Path

base = Path(r"F:\eeg_bci")

for p in [
    base / ".mne",
    base / "data" / "mne",
    base / "data" / "moabb",
    base / "tmp",
    base / "tmp" / "matplotlib",
]:
    p.mkdir(parents=True, exist_ok=True)

os.environ["_MNE_FAKE_HOME_DIR"] = str(base)
os.environ["MNE_DATA"] = str(base / "data" / "mne")
os.environ["MNE_DATASETS_BNCI_PATH"] = str(base / "data" / "mne")
os.environ["MOABB_DATA"] = str(base / "data" / "moabb")
os.environ["MPLCONFIGDIR"] = str(base / "tmp" / "matplotlib")
os.environ["JOBLIB_TEMP_FOLDER"] = str(base / "tmp")
\end{lstlisting}

然后在导入数据集前加入 MOABB Windows 路径修复：

\begin{lstlisting}[language=Python,caption={修复 MOABB 1.5 的 Windows 路径清洗问题}]
from pathlib import Path
import moabb.datasets.download as moabb_download

moabb_download._sanitize_path = lambda path: Path(path)

from moabb.datasets import BNCI2014_001
from moabb.paradigms import MotorImagery

dataset = BNCI2014_001()

print(dataset.event_id)
print(dataset.subject_list)
\end{lstlisting}

再读取第 1 个被试的左右手运动想象数据：

\begin{lstlisting}[language=Python,caption={读取 subject 1 的左右手运动想象数据}]
paradigm = MotorImagery(
    n_classes=2,
    events=["left_hand", "right_hand"],
    fmin=8,
    fmax=30,
    tmin=0.5,
    tmax=4.0,
    resample=128
)

X, y, metadata = paradigm.get_data(
    dataset=dataset,
    subjects=[1]
)

print("X shape:", X.shape)
print("y shape:", y.shape)
print("标签:", set(y))
print(metadata.head())
\end{lstlisting}

如果输出类似：

\begin{lstlisting}[language={}]
X shape: (288, 22, 449)
y shape: (288,)
标签: {'right_hand', 'left_hand'}
   subject session run
0        1  0train   0
1        1  0train   0
\end{lstlisting}

说明数据读取和预处理已经成功。这里 \texttt{288} 表示 trial 数量，\texttt{22} 表示 EEG 通道数量，\texttt{449} 表示每个 trial 的时间采样点数。设置 \texttt{tmin=0.5}、\texttt{tmax=4.0}、\texttt{resample=128} 后，理论时间长度为 \texttt{3.5 s}，MNE 通常会同时包含起止采样点，因此得到 \texttt{449} 个时间点是正常的。

需要注意：这一步只是下载、读取和预处理数据，还没有训练分类模型。真正训练发生在后续的 \texttt{cross\_val\_score} 或 \texttt{clf.fit(...)} 中。

\subsection{下载进度很慢和 SHA256 提示}

第一次读取公开数据集时，可能出现类似输出：

\begin{lstlisting}[language={}]
SHA256 hash of downloaded file: ...
Use this value as the 'known_hash' argument ...
11% | 4.64M/42.8M ...
\end{lstlisting}

\texttt{SHA256} 是文件校验值，用于确认下载文件是否完整，不是错误。下载速度慢通常是公开数据服务器或当前网络连接较慢导致的。只要最后能输出 \texttt{X shape}、\texttt{y shape} 和标签集合，就说明数据已经读取成功。下次运行同一数据通常会优先使用 F 盘缓存，不需要重新下载。

\subsection{确认数据没有主要写入 C 盘}

修复后，主要数据和缓存应该位于：

\begin{lstlisting}[language={}]
F:\eeg_bci\data\mne
F:\eeg_bci\data\moabb
F:\eeg_bci\.mne
F:\eeg_bci\tmp
\end{lstlisting}

可以在 Notebook 中检查：

\begin{lstlisting}[language=Python,caption={检查数据和配置路径}]
import os

print("MNE_DATA:", os.environ.get("MNE_DATA"))
print("MOABB_DATA:", os.environ.get("MOABB_DATA"))
print("MNE fake home:", os.environ.get("_MNE_FAKE_HOME_DIR"))
print("JOBLIB_TEMP_FOLDER:", os.environ.get("JOBLIB_TEMP_FOLDER"))
print("MPLCONFIGDIR:", os.environ.get("MPLCONFIGDIR"))
\end{lstlisting}

这些路径都指向 F 盘时，公开 EEG 数据集不会大量占用 C 盘。C 盘仍可能有少量系统级配置、浏览器缓存或 Windows 临时文件；如果运行内存不足，Windows 还可能使用 C 盘页面文件。这不是 EEG 数据主动下载到 C 盘。

\subsection{没有 Anaconda，为什么还能使用 \texttt{conda}}

本机没有安装 Anaconda。当前使用的 \texttt{conda} 来自安装在 F 盘的 Miniforge：

\begin{lstlisting}[language={}]
F:\Miniforge3\Scripts\conda.exe
\end{lstlisting}

三者区别可以简单理解为：

\begin{itemize}
  \item Anaconda：体积较大的 Python/Conda 发行版，默认包含大量科学计算包。
  \item Miniconda：精简版 Conda。
  \item Miniforge：基于 conda-forge 社区源的精简 Conda，适合科研和开源工具链。
\end{itemize}

因此，即使系统中直接输入 \texttt{conda} 找不到命令，也可以通过完整路径使用：

\begin{lstlisting}[language={},caption={使用 Miniforge 中的 conda}]
F:\Miniforge3\Scripts\conda.exe --version
\end{lstlisting}

本实验环境的位置是：

\begin{lstlisting}[language={}]
F:\eeg_bci\envs\eeg-mi
\end{lstlisting}

平时启动 JupyterLab 时也使用这个 Miniforge 的 \texttt{conda.exe}：

\begin{lstlisting}[language={},caption={启动 JupyterLab}]
F:\Miniforge3\Scripts\conda.exe run -p F:\eeg_bci\envs\eeg-mi jupyter lab --notebook-dir F:\eeg_bci\projects
\end{lstlisting}

\section{下肢脑电解码入门路线}

如果目标是做下肢相关脑电解码，需要先确认数据集中是否有下肢标签。当前入门使用的 \texttt{BNCI2014\_001} 是真实公开 EEG 数据集，不是仿真数据。它的事件标签为：

\begin{lstlisting}[language={}]
{'left_hand': 1, 'right_hand': 2, 'feet': 3, 'tongue': 4}
\end{lstlisting}

其中 \texttt{feet} 表示脚部或双脚运动想象类别，并不是 \texttt{left\_leg}、\texttt{right\_leg}、\texttt{left\_foot}、\texttt{right\_foot}。因此该数据集适合先做：

\begin{itemize}
  \item 左手 vs 右手。
  \item 左手 vs feet。
  \item 右手 vs feet。
  \item 手部 vs 脚部。
  \item 左手、右手、feet、tongue 四分类。
\end{itemize}

但它不能直接做：

\begin{lstlisting}[language={}]
左腿 vs 右腿
\end{lstlisting}

因为原始标签中没有左右腿或左右脚的区分。

\subsection{\texttt{BNCI2014\_001} 数据特点}

\texttt{BNCI2014\_001} 是经典运动想象 EEG 数据集，适合作为 BCI 入门和算法基线实验。它的主要特点如下：

\begin{itemize}
  \item 数据来自真实 EEG 采集，不是仿真数据。
  \item MOABB 中包含 9 名被试。
  \item 任务标签包括 \texttt{left\_hand}、\texttt{right\_hand}、\texttt{feet}、\texttt{tongue}。
  \item 常用 EEG 通道数为 22 个。
  \item 适合做运动想象分类、CSP + LDA 基线、频段分析和 ERD/ERS 分析。
\end{itemize}

在当前 Notebook 中使用：

\begin{lstlisting}[language=Python]
X, y, metadata = paradigm.get_data(
    dataset=dataset,
    subjects=[1]
)
\end{lstlisting}

得到的 \texttt{X} 是已经按任务事件切分好的 EEG 片段。若输出为：

\begin{lstlisting}[language={}]
X shape: (288, 22, 449)
y shape: (288,)
\end{lstlisting}

表示共有 288 个 trial，每个 trial 有 22 个 EEG 通道，每个通道有 449 个时间采样点。这里的 449 来自截取 \texttt{0.5--4.0 s} 并重采样到 128 Hz 后的时间点数；MNE 通常会同时包含起止采样点，因此比 \(3.5 \times 128 = 448\) 多 1 个点。

\subsection{\texttt{feet} 类别的含义}

\texttt{feet} 是脚部或双脚运动想象类别。它表示该 trial 中被试按照实验提示进行脚部相关运动想象，但数据标签没有进一步区分左脚、右脚、左腿或右腿。

因此，使用 \texttt{BNCI2014\_001} 时可以研究：

\begin{lstlisting}[language={}]
feet vs left_hand
feet vs right_hand
feet vs tongue
left_hand vs right_hand vs feet vs tongue
\end{lstlisting}

但不能研究：

\begin{lstlisting}[language={}]
left_foot vs right_foot
left_leg vs right_leg
\end{lstlisting}

原因是监督学习需要真实标签。若原始数据只有 \texttt{feet} 一个下肢标签，就不能通过代码把它拆成 \texttt{left\_foot} 和 \texttt{right\_foot}。强行拆分会造成伪标签，分类结果没有实验意义。

\subsection{用 \texttt{feet} 做什么分类}

当前阶段建议先做二分类，例如：

\begin{lstlisting}[language=Python,caption={提取 \texttt{left\_hand} vs \texttt{feet} 数据}]
paradigm = MotorImagery(
    n_classes=2,
    events=["left_hand", "feet"],
    fmin=8,
    fmax=30,
    tmin=0.5,
    tmax=4.0,
    resample=128
)
\end{lstlisting}

这个任务的含义是判断一个 EEG trial 更像左手运动想象，还是更像脚部运动想象。它不是最终的脊髓损伤左右腿解码任务，但可以作为下肢相关解码的入门基线。

也可以做：

\begin{lstlisting}[language=Python,caption={提取 \texttt{right\_hand} vs \texttt{feet} 数据}]
paradigm = MotorImagery(
    n_classes=2,
    events=["right_hand", "feet"],
    fmin=8,
    fmax=30,
    tmin=0.5,
    tmax=4.0,
    resample=128
)
\end{lstlisting}

如果后续目标是 SCI 下肢康复，\texttt{feet} 可以作为早期过渡任务，用于学习脑电处理流程、建立 CSP + LDA 基线、观察脚部运动想象在 Cz 附近的 Mu/Beta 节律变化。真正的左右腿或步态想象研究，需要进一步换用带有 \texttt{left\_leg/right\_leg} 或 \texttt{gait\_imagery/rest} 标签的数据集。

\subsection{为什么左右腿更难}

从脑电原理上看，左右手运动想象相对容易区分，因为左右手对应的感觉运动皮层分别位于对侧半球，空间差异较明显。下肢运动区更靠近大脑中线，左右腿在头皮 EEG 上的空间差异更小，而 EEG 本身空间分辨率有限，所以左右腿或左右脚解码通常比左右手更难。

\subsection{建议的下肢研究过渡方案}

如果当前阶段主要想做下肢脑电解码，建议先从 \texttt{feet} 类别开始，而不是直接追求左右腿分类。

可以先做二分类：

\begin{lstlisting}[language=Python,caption={左手 vs feet 下肢相关二分类}]
paradigm = MotorImagery(
    n_classes=2,
    events=["left_hand", "feet"],
    fmin=8,
    fmax=30,
    tmin=0.5,
    tmax=4.0,
    resample=128
)

X, y, metadata = paradigm.get_data(
    dataset=dataset,
    subjects=[1]
)

print("X shape:", X.shape)
print("y shape:", y.shape)
print("标签:", set(y))
\end{lstlisting}

也可以做右手 vs feet：

\begin{lstlisting}[language=Python,caption={右手 vs feet 下肢相关二分类}]
paradigm = MotorImagery(
    n_classes=2,
    events=["right_hand", "feet"],
    fmin=8,
    fmax=30,
    tmin=0.5,
    tmax=4.0,
    resample=128
)
\end{lstlisting}

在该阶段，分类器仍然可以使用 CSP + LDA。区别只是标签从 \texttt{left\_hand/right\_hand} 换成了 \texttt{left\_hand/feet} 或 \texttt{right\_hand/feet}。

\section{CSP + LDA 与 EEG 数据的数学耦合}

本节说明 CSP + LDA 到底在训练什么，以及它如何与当前提取到的 EEG 数据 \texttt{X} 和标签 \texttt{y} 连接。

\subsection{EEG 数据的数学形式}

通过 MOABB 得到的数据：

\begin{lstlisting}[language=Python]
X, y, metadata = paradigm.get_data(...)
\end{lstlisting}

可以看成一个三维数组：

\[
X \in \mathbb{R}^{N \times C \times T}
\]

其中：

\begin{itemize}
  \item \(N\) 是 trial 数量，例如 288。
  \item \(C\) 是 EEG 通道数，例如 22。
  \item \(T\) 是每个 trial 的时间采样点数，例如 449。
\end{itemize}

第 \(i\) 个 trial 可以写成矩阵：

\[
X_i \in \mathbb{R}^{C \times T}
\]

对应标签为：

\[
y_i \in \{0, 1\}
\]

例如在 \texttt{left\_hand vs feet} 任务中，可以把 \texttt{left\_hand} 视为类别 0，把 \texttt{feet} 视为类别 1。监督学习的目标是学习一个函数：

\[
f(X_i) \rightarrow y_i
\]

即输入一段 EEG trial，输出该 trial 属于哪一类运动想象。

\subsection{为什么先滤波}

运动想象 EEG 主要关注感觉运动节律，包括：

\[
\mu: 8\text{--}13~\mathrm{Hz}
\]

\[
\beta: 13\text{--}30~\mathrm{Hz}
\]

因此在代码中设置：

\begin{lstlisting}[language=Python]
fmin=8
fmax=30
\end{lstlisting}

相当于只保留与运动想象更相关的频段，减少低频漂移和高频噪声。滤波后，CSP 主要在 Mu/Beta 节律上寻找空间差异。

\subsection{CSP 的核心思想}

CSP 的目标是学习空间滤波器：

\[
W \in \mathbb{R}^{K \times C}
\]

其中 \(C\) 是原始 EEG 通道数，\(K\) 是选取的 CSP 分量数，例如代码中的：

\begin{lstlisting}[language=Python]
CSP(n_components=6)
\end{lstlisting}

对每个 trial，CSP 会进行空间投影：

\[
Z_i = W X_i
\]

其中：

\[
X_i \in \mathbb{R}^{C \times T}, \quad Z_i \in \mathbb{R}^{K \times T}
\]

也就是说，CSP 把原来的 22 个 EEG 通道变成若干个空间模式。每个空间模式都是原始电极通道的加权组合。

直观地说，CSP 会寻找这样的通道组合：

\[
z(t) = w_1 x_1(t) + w_2 x_2(t) + \cdots + w_C x_C(t)
\]

使得某一类任务中该组合信号的方差较大，而另一类任务中方差较小。

\subsection{协方差矩阵}

对第 \(i\) 个 trial，先计算归一化协方差矩阵：

\[
\Sigma_i = \frac{X_i X_i^\top}{\operatorname{trace}(X_i X_i^\top)}
\]

然后分别对两个类别求平均协方差：

\[
\Sigma_0 = \frac{1}{N_0} \sum_{i:y_i=0} \Sigma_i
\]

\[
\Sigma_1 = \frac{1}{N_1} \sum_{i:y_i=1} \Sigma_i
\]

CSP 要找的空间滤波方向 \(w\) 满足：投影后在一类中的方差尽量大，在另一类中的方差尽量小。这个目标可以写成：

\[
J(w) = \frac{w^\top \Sigma_0 w}{w^\top \Sigma_1 w}
\]

最大化或最小化这个比值，就能找到区分两类任务的空间模式。

这个优化问题通常转化为广义特征值问题：

\[
\Sigma_0 w = \lambda \Sigma_1 w
\]

特征值大的方向对应类别 0 方差较大的空间模式，特征值小的方向对应类别 1 方差较大的空间模式。实际使用时会取两端的若干个空间滤波器组成 \(W\)。

\subsection{CSP 输出的特征}

CSP 得到投影信号：

\[
Z_i = W X_i
\]

之后通常取每个空间分量的方差，并取对数：

\[
v_{ik} = \log \left( \operatorname{var}(Z_{ik}) \right)
\]

于是每个 EEG trial 不再是一个 \(C \times T\) 的时间序列矩阵，而变成一个 \(K\) 维特征向量：

\[
v_i \in \mathbb{R}^{K}
\]

这就是代码中 \texttt{CSP(log=True)} 的含义：用对数方差作为分类特征。

\subsection{LDA 的核心思想}

LDA 接收 CSP 输出的特征向量 \(v_i\)，学习一个线性分类器：

\[
g(v) = a^\top v + b
\]

分类规则为：

\[
\hat{y} =
\begin{cases}
0, & g(v) < 0 \\
1, & g(v) \geq 0
\end{cases}
\]

LDA 的目标是找到一个投影方向，使两类样本的均值分得尽量远，同时类内离散程度尽量小。

若两类 CSP 特征的均值分别为：

\[
\mu_0, \quad \mu_1
\]

类内散度矩阵为：

\[
S_w = \sum_{i:y_i=0}(v_i-\mu_0)(v_i-\mu_0)^\top
    + \sum_{i:y_i=1}(v_i-\mu_1)(v_i-\mu_1)^\top
\]

则 LDA 的典型投影方向可写为：

\[
a = S_w^{-1}(\mu_1 - \mu_0)
\]

也就是说，LDA 学的是一个能把两类 CSP 特征尽量分开的线性边界。

\subsection{从 EEG 到分类结果的完整链路}

当前实验中，CSP + LDA 的完整链路可以写成：

\[
X_i \in \mathbb{R}^{C \times T}
\quad \xrightarrow{\text{CSP}}
\quad v_i \in \mathbb{R}^{K}
\quad \xrightarrow{\text{LDA}}
\quad \hat{y}_i
\]

对应到代码就是：

\begin{lstlisting}[language=Python]
clf = make_pipeline(
    CSP(n_components=6, reg=None, log=True, norm_trace=False),
    LinearDiscriminantAnalysis()
)
\end{lstlisting}

\texttt{make\_pipeline} 的意义是把 CSP 和 LDA 串成一个整体。训练时，CSP 先从训练集 EEG 中学习空间滤波器 \(W\)，再把 EEG 转换成特征；LDA 再从这些特征中学习分类边界。

\subsection{交叉验证中的训练过程}

运行：

\begin{lstlisting}[language=Python]
scores = cross_val_score(clf, X, y, cv=cv, scoring="balanced_accuracy")
\end{lstlisting}

时，每一折都会执行：

\begin{enumerate}[label=\arabic*.]
  \item 用训练集估计 CSP 空间滤波器 \(W\)。
  \item 用训练集的 CSP 特征训练 LDA 分类边界。
  \item 把测试集 EEG 输入同一个 CSP，得到测试特征。
  \item 用 LDA 对测试特征分类。
  \item 计算测试集准确率。
\end{enumerate}

重要的是：CSP 必须只在训练集上学习，不能先对全部数据拟合 CSP 再交叉验证。把 CSP 和 LDA 放进 \texttt{Pipeline} 中，就是为了避免这种数据泄漏。

\subsection{在 \texttt{feet} 任务中的解释}

如果任务是：

\begin{lstlisting}[language={}]
left_hand vs feet
\end{lstlisting}

那么 CSP 学到的是能够区分“左手运动想象”和“脚部运动想象”的空间模式。它可能更关注 C3/C4 与 Cz 附近感觉运动区之间的差异。

如果任务是：

\begin{lstlisting}[language={}]
right_hand vs feet
\end{lstlisting}

则 CSP 学到的是“右手运动想象”和“脚部运动想象”的空间差异。

需要强调的是：在 \texttt{BNCI2014\_001} 中，\texttt{feet} 只有一个类别。因此 CSP + LDA 只能学习“脚部类 vs 另一个类”的差异，不能学习左脚和右脚之间的差异。

\subsection{与 SCI 下肢课题的关系}

对于脊髓损伤下肢运动想象解码，CSP + LDA 可以作为第一条传统基线。它的意义不是证明患者真的完成了下肢运动，而是检测 EEG 中是否存在与下肢运动想象相关、可被分类器稳定利用的感觉运动节律差异。

在后续更接近 SCI 的实验中，可以把标签替换成：

\begin{lstlisting}[language={}]
gait_imagery vs rest
left_leg vs right_leg
attempted_movement vs rest
\end{lstlisting}

但基本数学链路仍然相同：

\[
\text{EEG trial} \rightarrow \text{空间特征} \rightarrow \text{分类器} \rightarrow \text{运动意图标签}
\]

\subsection{如果必须做左右腿}

如果研究问题明确要求 \texttt{左腿 vs 右腿}，需要寻找带有左右腿、左右脚、踝关节或下肢单侧运动标签的数据集，或者自己设计采集实验。此时要重点确认：

\begin{itemize}
  \item 数据标签是否明确区分左腿和右腿。
  \item 是否是运动想象、实际运动，还是运动执行。
  \item 是否同步记录 EMG 或运动学信号，用于排除肌电和动作伪迹。
  \item 电极布局是否覆盖 Cz、FCz、CPz 及周围感觉运动区域。
  \item 试次数量是否足够进行交叉验证。
\end{itemize}

在没有左右腿标签的数据集中，不能通过代码把 \texttt{feet} 强行拆成左腿和右腿；那样没有真实标签支撑，分类结果没有实验意义。

\section{下一步建议}

建议下一步新建一个 Notebook，文件名可以为：

\begin{lstlisting}[language={}]
F:\eeg_bci\projects\motor_imagery_baseline.ipynb
\end{lstlisting}

先完成 \texttt{BNCI2014\_001} 的数据加载、事件查看和一个 CSP + LDA 二分类基线。这样后续无论是写课程项目、论文复现，还是接真实 EEG 设备，都会更清楚每一步在做什么。

\end{document}
